% 第 11 章 用auto\&\&完美传递（由 tools/convert.py 从原 content/ 自动生成，勿手改）

讨论了泛型代码中移动语义、完美转发参数后，现在讨论一下如何完美处理返回值。本章中，将讨论返回值的完美传递。为此，引入auto\&\&作为通用引用(同样也称为转发引用)的另一种方式。但也会讨论auto\&\&与转发值无关的部分。\par

下一章将讨论如何完美地返回值。\par


\section{默认的完美的传递}

我们经常需要将返回值传递给另一个函数:\par

\begin{codebox}[cpp]
// pass return value of compute() to process():
process(compute(t)); // OK, uses perfect forwarding of returned value
\end{codebox}

非泛型代码中，需要知道所涉及的类型。然而，泛型代码中，也希望\textit{compute()}的返回值完全传递给\textit{process()}。\par

好消息是：如果直接将返回值传递给另一个函数，该值会完美传递，保持其类型和值类别。不必担心移动语义(如果支持将自动使用)。\par

\mySmallsection{11.1.1 默认完美传递的细节}

完整的例子:\par

\begin{codeboxt}[cpp]{generic/perfectpassing.cpp}
#include <iostream>
#include <string>

void process(const std::string&) {
  std::cout << "process(const std::string&)\n";
}
void process(std::string&) {
  std::cout << "process(std::string&)\n";
}
void process(std::string&&) {
  std::cout << "process(std::string&&)\n";
}

const std::string& computeConstLRef(const std::string& str) {
  return str;
}
  std::string& computeLRef(std::string& str) {
  return str;
}
  std::string&& computeRRef(std::string&& str) {
  return std::move(str);
}
  std::string computeValue(const std::string& str) {
  return str;
}

int main()
{
  process(computeConstLRef("tmp")); // calls process(const std::string&)

  std::string str{"lvalue"};
  process(computeLRef(str)); // calls process(std::string&)

  process(computeRRef("tmp")); // calls process(std::string&&)
  process(computeRRef(std::move(str))); // calls process(std::string&&)

  process(computeValue("tmp")); // calls process(std::string&&)
}
\end{codeboxt}

\begin{itemize}
  \item 如果\textit{compute()}返回一个const lvalue引用:\par
\begin{codebox}[cpp]
  const std::string& computeConstLRef(const std::string& str) {
    return str;
  }
\end{codebox}
  返回值的值类别是lvalue，这意味着返回值会完美转发，并与\textit{const} lvalue引用匹配:\par
\begin{codebox}[cpp]
  process(computeConstLRef("tmp")); // calls process(const std::string&)
\end{codebox}
  \item 如果\textit{compute()}返回一个非\textit{const} lvalue引用:\par
\begin{codebox}[cpp]
  std::string& computeLRef(std::string& str) {
    return str;
  }
\end{codebox}
  返回值的值类别是lvalue，这意味着返回值会完全转发，并与非\textit{const} lvalue引用的最佳匹配:\par
\begin{codebox}[cpp]
  std::string str{"lvalue"};
  process(computeLRef(str)); // calls process(std::string&)
\end{codebox}
  \item 如果\textit{compute()}返回rvalue引用:\par
\begin{codebox}[cpp]
  std::string&& computeRRef(std::string&& str) {
    return std::move(str);
  }
\end{codebox}
  返回值的值类别是xvalue，这意味着返回值会完全转发为rvalue引用，允许\textit{process()}窃取值:\par
\begin{codebox}[cpp]
  process(computeRRef("tmp")); // calls process(std::string&&)
  process(computeRRef(std::move(str))); // calls process(std::string&&)
\end{codebox}
  \item 如果\textit{compute()}按值返回临时对象:\par
\begin{codebox}[cpp]
  std::string computeValue(const std::string& str) {
    return str;
  }
\end{codebox}
  返回值的值类别是prvalue，返回值完全转发为rvalue引用，也允许\textit{process()}窃取值:\par
\begin{codebox}[cpp]
  process(computeValue("tmp")); // calls process(std::string&&)
\end{codebox}
\end{itemize}

注意，通过返回\textit{const}值:\par

\begin{codebox}[cpp]
const std::string computeConstValue(const std::string& str) {
  return str;
}
\end{codebox}

或\textit{const} rvalue引用:\par

\begin{codebox}[cpp]
const std::string&& computeConstRRef(std::string&& str) {
  return std::move(str);
}
\end{codebox}

再次禁用移动语义:\par

\begin{codebox}[cpp]
process(computeConstValue("tmp")); // calls process(const std::string&)
process(computeConstRRef("tmp")); // calls process(const std::string&)
\end{codebox}

如果有\textit{const}\&\&的声明，可以接受这样的重载。\par

因此:不要将value返回的值标记为\textit{const}，也不要将返回的非\textit{const} rvalue引用标记为\textit{const}。\par


\section{使用auto\&\&的通用引用}

但在泛型代码中，如何在传递返回值的同时，仍然保持其类型和值类别?\par

答案是通用/转发引用，但没有声明为参数。为此，需要auto\&\&。\par

调用时:\par

\begin{codebox}[cpp]
// pass return value of compute() to process():
process(compute(t)); // OK, uses perfect forwarding of returned value
\end{codebox}

还可以实现以下功能:\par

\begin{codebox}[cpp]
// pass return value of compute() to process() with some delay:
auto&& ret = compute(t); // initialize a universal reference with the return value
...
process(std::forward<decltype(ret)>(ret)); // OK, uses perfect forwarding of returned value
\end{codebox}

或者，使用大括号初始化时:\par

\begin{codebox}[cpp]
// pass return value of compute() to process() with some delay:
auto&& ret{compute(t)}; // initialize a universal reference with the return value
...
process(std::forward<decltype(ret)>(ret)); // OK, uses perfect forwarding of returned value
\end{codebox}

请参阅generic/perfectautorefref.cpp以获得所有情况下的完整示例。\par

\mySmallsection{11.2.1 auto\&\&的类型定义}

使用auto\&\&声明时，也声明了一个通用引用。定义一个绑定到所有值类别的引用，该引用的类型保留其初始值的类型和值类别。\par

如果声明\par

\begin{codebox}[cpp]
auto&& ref{ ... }; // ref is a universal/forwarding reference
\end{codebox}

类型ref是根据函数模板参数的通用引用类型推导出来的:\par

\begin{codebox}[cpp]
template<typename T>
void callFoo(T&& ref); // ref is a universal/forwarding reference
\end{codebox}

根据规则，ref(即auto\&\&或T\&\&类型)的类型是\par

\begin{itemize}
  \item 如果引用lvalue，则为lvalue引用
  \item 如果引用rvalue，则为rvalue引用
\end{itemize}

例如:\par

\begin{codebox}[cpp]
// forward declarations:
std::string retByValue();
std::string& retByRef();
std::string&& retByRefRef();
const std::string& retByConstRef();
const std::string&& retByConstRefRef();

// deduced auto&& types:
std::string s;
auto&& r1{s}; // std::string&
auto&& r2{std::move(s)}; // std::string&&

auto&& r3{retByValue()}; // std::string&&
auto&& r4{retByRef()}; // std::string&
auto&& r5{retByRefRef()}; // std::string&&
auto&& r6{retByConstRef()}; // const std::string&
auto&& r7{retByConstRefRef()}; // const std::string&&
\end{codebox}

当使用rvalue(prvalue或xvalue)初始化用auto\&\&声明的引用时，就声明了一个rvalue引用。例如，将引用绑定到标记为\textit{std::move()}的对象或返回的普通值或rvalue引用时，就会出现这种情况。\par

但是，当使用lvalue初始化用auto\&\&声明的引用时，就声明了一个lvalue引用。例如，将引用绑定到已命名的对象或返回lvalue引用的函数的返回值时，就会出现这种情况。\par

因为字符串字面值是lvalue(以字符数组为类型)，所以将通用引用绑定到字符串字面值时，也会得到lvalue引用:\par

\begin{codebox}[cpp]
auto&& r8{"hello"}; // const char(&)[6]
\end{codebox}

因为对函数的引用总是lvalue，所以将通用引用绑定到函数时，也会得到lvalue引用:\par

\begin{codebox}[cpp]
std::string foo(int); // forward declaration

auto&& r9{foo}; // lvalue of type std::string(&)(int)
\end{codebox}

\mySmallsection{11.2.2 完美转发auto\&\&引用}

同样，也可以完美地将传递给通用引用的值作为函数模板形参:\par

\begin{codebox}[cpp]
template<typename T>
void callFoo(T&& ref) {
  foo(std::forward<T>(ref)); // becomes foo(std::move(ref)) for passed rvalues
}
\end{codebox}

完全可以转发auto\&\&的通用引用:\par

\begin{codebox}[cpp]
auto&& ref{ ... };

foo(std::forward<decltype(ref)>(ref)); // becomes foo(std::move(ref)) for rvalues
\end{codebox}

表达式\textit{std::\allowbreak{}forward<\allowbreak{}decltype(\allowbreak{}ref)\allowbreak{}>(\allowbreak{}ref)\allowbreak{}}是这样实现的:\par

\begin{itemize}
  \item 如果传递的decltype(ref)类型是一个lvalue引用，如果ref是用返回的lvalue引用初始化的，则表达式将ref强制转换为lvalue引用，这意味着ref传递时没有移动语义。
  \item 如果传递的decltype(ref)类型是rvalue引用，如果ref是用返回的普通值或rvalue引用初始化的，则表达式将ref转换为rvalue引用，这是\textit{std::\allowbreak{}move(\allowbreak{}ref)\allowbreak{}}的效果。
\end{itemize}

因此，如果用函数的返回值初始化通用引用:\par

\begin{codebox}[cpp]
auto&& ret{compute(t)}; // initialize a universal reference with the return value
\end{codebox}

表达式为\par

\begin{codebox}[cpp]
process(std::forward<decltype(ret)>(ret)); // perfectly forward the return value
\end{codebox}

当\textit{compute()}返回rvalue(如临时对象或rvalue引用)时，扩展为\par

\begin{codebox}[cpp]
process(std::move(ret));
\end{codebox}


\section{auto\&\&的非转发引用}

再次注意，通用引用是将引用绑定到任何类型和值类别的任何对象，并且仍然保留其值类别和是否为\textit{const}的唯一方法。这也适用于用auto\&\&声明的通用引用。\par

\mySmallsection{11.3.1 通用引用和基于范围的for循环}

使用基于范围的for循环时，使用auto\&\&声明的非转发通用引用起着重要的作用。\par

\mySmallsection{基于范围的循环规范}

在C++标准中，基于范围的for循环可以指定使用普通for循环遍历范围内的元素。\par

类似这样的调用:\par

\begin{codebox}[cpp]
std::vector<std::string> coll;
...
for (const auto& s : coll) {
  ...
}
\end{codebox}

相当于:\par

\begin{codebox}[cpp]
std::vector<std::string> coll;
...
auto&& range = coll; // initialize a universal reference
auto pos = range.begin(); // to use the given range coll here
auto end = range.end(); // and here
for ( ; pos != end; ++pos ) {
  const auto& s = *pos;
  ...
}
\end{codebox}

将range声明为一个通用引用，希望能够将其绑定到每个range，所以可以使用它两次(一次是开始，一次是结束)没有创建副本或丢失信息。\par

循环应该为:\par

\begin{itemize}
  \item 非\textit{const} lvalue:
\begin{codebox}[cpp]
  std::vector<int> coll;
  ...
  for (int& i : coll) {
    i *= 2;
  }
\end{codebox}
  \item \textit{const} lvalue:
\begin{codebox}[cpp]
  const std::vector<int> coll{0, 8, 15};
  ...
  for (int i : coll) {
    ...
  }
\end{codebox}
  \item prvalue:
\begin{codebox}[cpp]
  for (int i : std::vector<int>{0, 8, 15}) {
    ...
  }
\end{codebox}
\end{itemize}

注意，对于这些情况，没有其他方法声明range:\par

\begin{itemize}
  \item 使用auto，将创建range的副本(这需要花费时间并禁用对元素的修改)。
  \item 使用auto\&，可以用临时的prvalue禁用range的初始化。
  \item 使用const auto\&，将失去所遍历range的非常量性。
\end{itemize}

请注意，现在指定的基于范围的for循环有一个问题。代码如以下:\par

\begin{codebox}[cpp]
  std::vector<std::string> createStrings();
  ...
  for (char c : createStrings().at(0)) { // fatal runtime error
    ...
  }
\end{codebox}

成为:\par

\begin{codebox}[cpp]
  std::vector<std::string> createStrings();
  ...
  auto&& range = createStrings().at(0); // OOPS: universal reference to reference
  auto pos = range.begin(); // return value of createStrings() destroyed here
  auto end = range.end();
  for ( ; pos != end; ++pos ) {
    char c = *pos;
    ...
  }
\end{codebox}

所有有效引用都延长了所绑定值的生命周期，这也适用于rvalue引用。但是，没有绑定到\textit{createString(\allowbreak{})\allowbreak{}}的返回值(这样可以正常工作);而是绑定到引用，该引用指向由\textit{at()}返回的\textit{createStrings(\allowbreak{})\allowbreak{}}的返回类型，扩展了引用的生命周期。因此，该循环将遍历已经销毁的字符串。\par

\mySmallsection{使用基于范围的for循环}

即使在调用基于范围的for循环时，通用引用也有意义。\par

要在迭代时修改元素，必须使用非\textit{const}引用。考虑函数模板，将传递值赋给传递集合中的所有元素:\par

\begin{codebox}[cpp]
template<typename Coll, typename T>
void assign(Coll& coll, const T& value) {
  for (auto& elem : coll) {
    elem = value;
  }
}
\end{codebox}

看起来适用于所有容器类型和元素类型(其中支持赋值):\par

\begin{codebox}[cpp]
std::vector<int> coll1{0, 8, 15};
...
assign(coll1, 42); // OK

std::vector<std::string> coll2{"hello", "world"};
...
assign(coll2, "ok"); // OK
\end{codebox}

然而，有种情况行不通:\par

\begin{codebox}[cpp]
std::vector<bool> collB{false, true, false};
...
assign(collB, true); // ERROR: cannot bind non-const lvalue reference to an rvalue
\end{codebox}

发生了什么事?看一下基于范围的for循环展开的代码:\par

\begin{codebox}[cpp]
std::vector<bool> coll{false, true, false};
...
{
  auto&& range = coll; // OK: universal reference to reference
  auto pos = range.begin(); // OK
  auto end = range.end(); // OK
  for ( ; pos != end; ++pos ) { // OK
    auto& elem = *pos; // ERROR: cannot bind non-const lvalue reference to an rvalue
    elem = elem + elem;
  }
}
\end{codebox}

问题是std::vector<bool>中的元素不是bool类型的对象，而是单个比特位。实现方法是:对于std::vector<bool>，元素引用的类型不是元素类型的引用。std::\allowbreak{}vector<\allowbreak{}bool>的实现是主模板std::vector<T>实现的偏特化，其中元素的引用是代理类的对象，可以像引用一样使用:\par

\begin{codebox}[cpp]
namespace std {
  template< ... >
  class vector<bool, ... > {
    public:
    ...
    class reference {
      ...
    };
    ...
  };
}
\end{codebox}

当对迭代器解引用时，返回std::vector<bool>::reference的值。因此，在基于范围的for循环的扩展代码中的语句为\par

\begin{codebox}[cpp]
auto& elem = *pos;
\end{codebox}

尝试将非\textit{const} lvalue引用绑定到临时对象(prvalue)，这是不允许的。\par

然而，这个问题有一个解决方案:调用基于范围的for循环时使用通用引用:\par

\begin{codebox}[cpp]
template<typename Coll, typename T>
void assign(Coll& coll, const T& value) {
  for (auto&& elem : coll) { // note: universal reference support proxy element types
    elem = value;
  }
}
\end{codebox}

因为通用引用可以绑定到任何对象(甚至是prvalue)，所以vector<bool>的代码现在可以编译:\par

\begin{codebox}[cpp]
std::vector<bool> collB{false, true, false};
...
assign(collB, true); // OK (universal reference used to bind to an element)
\end{codebox}

因此，找到了使用非转发通用引用的另一个原因:绑定到没有作为引用实现的引用类型。或者说:允许绑定作为代理类型提供的非\textit{const}对象来进行操作。\par


\section{完美的Lambda转发}

如果想要完美地转发通用Lambda的参数，必须使用通用引用和\textit{std::\allowbreak{}forward<\allowbreak{}>(\allowbreak{})\allowbreak{}}。但现在只需使用auto\&\&来声明通用引用。\par

已经有了一个函数模板:\par

\begin{codebox}[cpp]
auto callFoo = [](auto&& arg) { // arg is a universal/forwarding reference
  foo(std::forward<decltype(arg)>(arg)); // perfectly forward arg
};

std::string s{"BLM"};
callFoo(s); // OK, arg is std::string&
callFoo(std::move(s)); // OK, arg is std::string&&
\end{codebox}

在C++20中，可以通过用模板形参声明Lambda来避免使用decltype(arg)。\par

下面的泛型Lambda使用此函数，完美地转发了可变数量的参数:\par

\begin{codebox}[cpp]
[] (auto&&... args) {
  ...
  foo(std::forward<decltype(args)>(args)...);
};
\end{codebox}

请记住Lambda只是定义函数对象的一种简单方法(定义了operator()的对象允许其作为函数使用)。上面的定义扩展为编译器定义的类(闭包类型)，并将通用引用定义为模板形参:\par

\begin{codebox}[cpp]
class NameDefinedByCompiler {
  ...
  public:
  template<typename... Args>
  auto operator() (Args&&... args) const {
    ...
    foo(std::forward<decltype(args)>(args)...);
  }
};
\end{codebox}

再次注意，用\textit{const}(或\textit{volatile})限定的泛型rvalue引用不是通用引用，这也适用于Lambda。如果形参用\textit{const} auto\&\&声明，只能传递rvalue:\par

\begin{codebox}[cpp]
auto callFoo = [](const auto&& arg) { // arg is not a universal reference
  ...
};

const std::string cs{"BLM"};
callFoo(cs); // ERROR: s is not an rvalue
callFoo(std::move(cs)); // OK, arg is const std::string&&
\end{codebox}

Lambda内部，使用\textit{std::move()}完美传递传递参数就可以了。\par


\section{C++20中函数声明使用auto\&\&}

因为C++20，也可以用auto\&\&来声明普通函数，处理方式一样:用通用引用来声明函数模板。\par

如下定义:\par

\begin{codebox}[cpp]
void callFoo(auto&& val) {
  foo(std::forward<decltype(arg)>(arg));
}
\end{codebox}

相当于:\par

\begin{codebox}[cpp]
template<typename T>
void callFoo(T&& val) {
  foo(std::forward<decltype(arg)>(arg));
}
\end{codebox}

还可以在这里显式地指定参数的类型。然而，与普通函数模板一样，必须用适合所传递参数的值类别的类型来限定模板形参:\par

\begin{codebox}[cpp]
std::string s;
...
callFoo<std::string>(s); // ERROR: cannot bind rvalue reference to lvalue
callFoo<std::string&>(s); // OK, does not move and forward s
callFoo<std::string>(std::move(s)); // OK, does move and forward s
callFoo<std::string&&>(std::move(s)); // OK, does move and forward s
\end{codebox}


\section{总结}

\begin{itemize}
  \item 不要使用\textit{const}返回值(否则将禁用返回值的移动语义)。
  \item 不要将返回的非\textit{const} rvalue引用标记为\textit{const}。
  \item auto\&\&可用于声明不是参数的通用引用。就像任何通用引用一样，它可以引用任何类型和值类别的所有对象，其类型是
  \begin{itemize}
    \item[-] 如果绑定到lvalue，则为lvalue引用(类型\&)
    \item[-] 如果绑定到rvalue，则为rvalue引用(类型\&\&)
  \end{itemize}
  \item 使用\textit{std::\allowbreak{}forward<\allowbreak{}decltype(\allowbreak{}ref)\allowbreak{}>(\allowbreak{}ref)\allowbreak{}}完美地转发用auto\&\&声明的通用引用ref。
  \item 可以使用通用引用同时引用\textit{const}和非\textit{const}对象，并多次使用而不会丢失\textit{const}信息。
  \item 可以使用通用引用来绑定到代理类型的引用。
  \item 考虑在迭代集合元素以修改元素的泛型代码中使用auto\&\&。这样，代码就可以用于代理类型的引用。
  \item 声明Lambda(或函数，因为C++20)的形参时使用auto\&\&是在函数模板中声明作为通用引用形参的方式。
\end{itemize}
